\section{Introduction and claim ledger}
\label{sec:introduction}

The fixed-shift prime-pair problem asks for asymptotics of
\(
 \sum_n\Lambda(n)\Lambda(n+h)W(n/X)
\)
when $h$ is fixed.  For $h=2$ this contains the twin-prime conjecture.
Classical sieve and distribution theorems give powerful one-sided or
averaged information, but they do not supply the required fixed-shift
two-prime lower bound; see, for example,
\citep{HardyLittlewood1923,IwaniecKowalski2004,FriedlanderIwaniec1998,
FordMaynard2024}.  The purpose of this paper is narrower.  We study which
parts of a related factorized form can be localized using established
theorems, and we quantify exactly what is not determined by finitely many
Mellin coordinates.

Fix throughout an integer $h\ne0$ and a real-valued function
$W\in C_c^1((0,\infty))$.  Terms for which $mn+h\leq0$ are understood to
vanish.  Put
\begin{equation}\label{eq:intro-residual-measure}
 \nu_{X,h}
 =\sum_{m,n\geq1}
   (\Lambda(n)-1)\Lambda(mn+h)W(mn/X)
   \delta_{\log(n/m)}.
\end{equation}
Its Fourier--Mellin transform is the finite exponential polynomial
\begin{equation}\label{eq:intro-mellin-transform}
 \widehat\nu_{X,h}(t)
 =\int_{\R}e^{itu}\,d\nu_{X,h}(u)
 =\sum_{m,n\geq1}(\Lambda(n)-1)\Lambda(mn+h)W(mn/X)
   (n/m)^{it}.
\end{equation}
For $V\in C_c^1(\R)$ and a center $s\in\R$, the corresponding local
factor-scale observation is
\begin{equation}\label{eq:intro-local-window}
 \cA_{h,V}(X;s)
 =\int_{\R}V(u-s)\,d\nu_{X,h}(u).
\end{equation}

The preceding paper \citep{WangTPCNine2026} evaluated only the total mass
$\widehat\nu_{X,h}(0)$.  Write
\begin{equation}\label{eq:intro-Ij}
 I_0(W)=\int_0^\infty W(u)\,du,
 \qquad
 I_1(W)=\int_0^\infty W(u)\log u\,du,
\end{equation}
and
\begin{equation}\label{eq:intro-Chs}
 C_h(z)=
 \prod_{p\mid h}(1-p^{-z-1})
 \prod_{p\nmid h}\left(1+\frac1{(p-1)p^{z+1}}\right),
 \qquad C_h=C_h(0).
\end{equation}
The exact divisor collapse
\begin{equation}\label{eq:intro-zero-collapse}
 \sum_{d\mid r}(\Lambda(d)-1)=\log r-\tau(r)
\end{equation}
and the uniform Titchmarsh theorem of Assing, Blomer and Li
\citep{AssingBlomerLi2021} give, for every $A>0$,
\begin{align}
 \widehat\nu_{X,h}(0)
 =X\Big[&(1-C_h)\{(\log X)I_0(W)+I_1(W)\}
 \notag\\
 &-2\{\gamma C_h+C_h'(0)\}I_0(W)\Big]
 +O_{A,h,W}\!\left(X(\log X)^{-A}\right).
 \label{eq:intro-complete-zero-mode}
\end{align}
Complete factor aggregation is decisive in
\eqref{eq:intro-zero-collapse}; it also destroys the location of the
individual divisor $d$.  The present paper separates three increasingly
fine questions.

\subsection{The positive results}

The first result is a coarse but genuinely two-sided localization.  If
$m,n\leq Y$ cannot occur on the support of $W$---as is automatic for
$Y\leq X^{1/3}$ and large $X$---then removing the two hard edges costs at
most $O(X\log(2Y))$.  Consequently the leading $X\log X$ coefficient in
\eqref{eq:intro-complete-zero-mode} remains in the interior $m,n>Y$ when
$\log Y=o(\log X)$.  The same statement holds for the unsigned
prime-source form with coefficient $1$; see
\cref{thm:bulk-persistence}.  This does not confine $n/m$ to a fixed
dyadic window.  It says that the leading aggregate is not carried by any
fixed, polylogarithmic, or subpower collection of factor-edge layers.

The second result is fine localization on the directly accessible side.
For
\begin{equation}\label{eq:intro-Type-I-range}
 D\leq X^{1/2}(\log X)^{-B},
\end{equation}
maximal Bombieri--Vinogradov evaluates
\begin{equation}\label{eq:intro-Type-I-local-form}
 \sum_{\substack{n\leq D\\m\geq1}}
 (\Lambda(n)-1)\Lambda(mn+h)W(mn/X)
 V(\log(n/m)-s)
\end{equation}
with an explicit main term, uniformly for every real center $s$; see
\cref{thm:Type-I-scale-localization}.  Thus every fixed smooth log-scale
window is available in the Type~I range.  The proof is an organization
of the classical theorem, not an improvement of its distribution level.

\subsection{The information boundary}

A local window is an integral of the whole Mellin transform, rather than
a function of finitely many frequencies.  We make that distinction exact.
For a total-variation ball of coefficient measures, arbitrary nonlinear
reconstruction from finitely many samples and jets has minimax error
\begin{equation}\label{eq:intro-minimax-preview}
 R\,\dist_\infty\!\left(
 V,\operatorname{span}\{u^k e^{it_\ell u}:0\leq k<s_\ell\}
 \right).
\end{equation}
This is the standard quotient-space principle of optimal recovery applied
to the factor-scale coordinate; compare \citep{MicchelliRivlin1977}.
What is arithmetic here is the explicit realization of null directions on
actual divisor geometries.  The factor pairs of $r=p^D$ form an equally
spaced log-scale grid, so a polynomial with prescribed unit-circle roots
annihilates any specified finite set of samples and jets.

For $h=2$, the row $r=3^4=81$ has prime target $83$.  On its five factor
nodes we exhibit two probability profiles having the same Mellin jets
through order three but different central-window mass.  This is an exact
certificate on genuine divisor nodes and a genuine prime-target row, but
the two profiles are deliberately adversarial: they are not the arithmetic
weights $\Lambda(d)-1$.  The minimax theorem therefore rules out a
deduction from finite Mellin data plus a norm bound alone.  It does not
assert that the actual prime residual realizes a worst-case direction.

\subsection{The first arithmetic jets and the endpoint boundary}

For a fixed product $r$, define
\[
 B_r(t)=\sum_{d\mid r}(\Lambda(d)-1)
 e^{it(2\log d-\log r)}.
\]
The zero coefficient is \eqref{eq:intro-zero-collapse}.  The next two
coefficients already contain two-von-Mangoldt convolutions:
\begin{align*}
 i^{-1}B_r'(0)
  &=(\log r)^2-2(\Lambda*\Lambda*\one)(r),\\
 i^{-2}B_r''(0)
  &=(\log r)^3-Q_2(r)
    -4((\mathfrak D\Lambda)*\Lambda*\one)(r),
\end{align*}
where $(\mathfrak Df)(n)=f(n)\log n$ and $Q_2$ is explicit.  Even
symmetrization removes the odd coefficient, but not the second one.
These are exact identities, not cancellation estimates.

At the opposite extreme, an all-frequency kernel after a two-variable
Mellin demodulation selects the $m=1$ layer exactly.  Its evaluation with
the Hardy--Littlewood main term is therefore equivalent to the weighted
fixed-shift prime-pair asymptotic, rather than a proof of it.

\subsection{One-way claim ledger}

For clarity, the logical content is summarized in
\cref{tab:intro-claim-ledger}.

\begin{table}[ht]
\centering
\small
\begin{tabular}{@{}
 >{\raggedright\arraybackslash}p{0.25\textwidth}
 >{\raggedright\arraybackslash}p{0.31\textwidth}
 >{\raggedright\arraybackslash}p{0.34\textwidth}@{}}
\toprule
Input & Proved output & Not obtained \\
\midrule
Complete zero mode + upper sieves
& Bulk persistence after $m,n\leq Y$ are deleted
& Fixed-ratio or signed Type~II cancellation \\
Maximal Bombieri--Vinogradov
& Smooth scale localization for $n\leq D$
& Localization of the complete $n>D$ tail \\
Finite Mellin samples/jets + a TV bound
& Exact minimax radius and null profiles
& Impossibility for the actual $\Lambda-1$ profile \\
Taylor and Fourier identities
& Explicit parity-sensitive interfaces and exact $m=1$ selector
& Hardy--Littlewood asymptotic or twin-prime lower bound \\
\bottomrule
\end{tabular}
\caption{Every implication in the paper is one-way.}
\label{tab:intro-claim-ledger}
\end{table}
